\documentclass[12pt,a4paper]{article}

% ========= Codificación / Idioma =========
\usepackage[utf8]{inputenc} % seguro para pdfLaTeX
\usepackage[T1]{fontenc}
\usepackage[spanish]{babel}
\usepackage{lmodern}

% ========= Márgenes =========
\usepackage[a4paper,margin=2.5cm]{geometry}

% ========= Listas (NECESARIO para opciones en enumerate/itemize) =========
\usepackage{enumitem}

% ========= Gráficos / color =========
\usepackage{graphicx}
\usepackage{xcolor}
\usepackage{tikz}
\usetikzlibrary{calc}

% ========= Enlaces / índice =========
\usepackage{hyperref}
\hypersetup{
  colorlinks=true,
  linkcolor=black,
  urlcolor=blue
}

% ========= Encabezado / pie =========
\usepackage{fancyhdr}
\usepackage{lastpage}

% ========= Espaciado =========
\usepackage{setspace}

% ========= Colores similares a TecNM =========
\definecolor{tecnmOrange}{RGB}{243,110,33}
\definecolor{tecnmBlack}{RGB}{0,0,0}

% ========= Helper: logo con fallback =========
\newcommand{\LogoBox}[3]{%
  \IfFileExists{#1}{\includegraphics[height=#2]{#1}}{%
    \fbox{\parbox[c][#2][c]{3.6cm}{\centering\footnotesize #3}}%
  }%
}

% =========================
% DATOS EDITABLES
% =========================
\newcommand{\Carrera}{Ingeniería en Sistemas Computacionales}
\newcommand{\TipoDocumento}{REPORTE SEMANA 1 DE RESIDENCIA PROFESIONAL}
\newcommand{\TituloProyecto}{Diseño de kiosco digital para el proceso de credencialización en el TecNM Campus Oaxaca}
\newcommand{\Alumno}{Mendoza Luis Prometeo}
\newcommand{\NoControl}{21160727}
\newcommand{\Semestre}{10\textsuperscript{o}}
\newcommand{\Periodo}{Enero--Junio 2026}
% Evitamos el guion largo Unicode “—” para que no falle pdfLaTeX:
\newcommand{\LugarFecha}{Oaxaca de Juárez, Oaxaca --- Enero 2026}

% ====== Rutas robustas para logos (sin sufrir) ======
\newcommand{\LogoAuto}[3]{%
  % #1 = nombre archivo con extensión
  % #2 = altura objetivo
  % #3 = texto fallback
  \IfFileExists{../assets/logos/#1}{%
    \includegraphics[height=#2,width=\linewidth,keepaspectratio]{../assets/logos/#1}%
  }{%
    \IfFileExists{assets/logos/#1}{%
      \includegraphics[height=#2,width=\linewidth,keepaspectratio]{assets/logos/#1}%
    }{%
      \fbox{\parbox[c][#2][c]{3.6cm}{\centering\footnotesize #3}}%
    }%
  }%
}

\newlength{\HeaderShift}
\setlength{\HeaderShift}{2.0cm} % <-- sube/baja esto (2.0–2.6cm suele quedar bien)


\begin{document}

% =========================
% PORTADA (Página 1)
% =========================
\pagenumbering{arabic}
\setcounter{page}{1}
\thispagestyle{empty}

% DIBUJO DE LÍNEAS (como la portada)
\begin{tikzpicture}[remember picture, overlay]
  % Línea horizontal naranja (arriba)
  \draw[tecnmOrange, line width=1.2pt]
    ($(current page.north west)+(3.2cm,-3.6cm)$) --
    ($(current page.north east)+(-2.2cm,-3.6cm)$);

  % Barras verticales izquierda (negra + naranja + negra)
  \draw[tecnmBlack, line width=3.0pt]
    ($(current page.north west)+(2.0cm,-4.2cm)$) --
    ($(current page.south west)+(2.0cm,2.0cm)$);

  \draw[tecnmOrange, line width=2.0pt]
    ($(current page.north west)+(2.5cm,-4.2cm)$) --
    ($(current page.south west)+(2.5cm,2.0cm)$);

  \draw[tecnmBlack, line width=0.8pt]
    ($(current page.north west)+(2.9cm,-4.2cm)$) --
    ($(current page.south west)+(2.9cm,2.0cm)$);

  % Número de página abajo derecha (usa el contador real)
  \node[anchor=south east] at ($(current page.south east)+(-1.6cm,1.2cm)$) {\small \thepage};
\end{tikzpicture}

% CABECERA CON LOGOS
\vspace*{0.8cm}

\noindent
\hspace*{\HeaderShift}%
\begin{minipage}[t]{\dimexpr\textwidth-\HeaderShift\relax}

  \begin{minipage}[t]{0.34\linewidth}
    \raggedright
    \LogoAuto{SecretariaEducacion.png}{1.5cm}{LOGO EDUCACIÓN}
  \end{minipage}
  \hfill
  \begin{minipage}[t]{0.40\linewidth}
    \centering
    \LogoAuto{Logo-TecNM.png}{1.6cm}{LOGO TECNM}
  \end{minipage}
  \hfill
  \begin{minipage}[t]{0.20\linewidth}
    \raggedleft
    \LogoAuto{Instituto_Tecnologico_de_Oaxaca.png}{1.6cm}{LOGO ITO}
  \end{minipage}

\end{minipage}

\vspace{2.2cm}


% TEXTO CENTRAL
\begin{center}
  \setstretch{1.15}

  {\large \textbf{INSTITUTO TECNOLÓGICO DE OAXACA}\par}
  \vspace{0.4cm}
  {\large \textbf{DIVISIÓN DE ESTUDIOS PROFESIONALES}\par}

  \vspace{1.2cm}
  {\normalsize \textbf{\Carrera}\par}

  \vspace{1.2cm}
  {\normalsize \textbf{\TipoDocumento}\par}

  \vspace{1.2cm}
  {\normalsize \TituloProyecto\par}

  \vspace{1.6cm}
  {\normalsize \textbf{PRESENTA:}\par}
  \vspace{0.2cm}
  {\normalsize \textbf{\Alumno}\par}

  \vspace{1.6cm}
  {\normalsize \textbf{NÚMERO DE CONTROL:}\par}
  \vspace{0.2cm}
  {\normalsize \textbf{\NoControl}\par}

  \vspace{2.0cm}
  {\normalsize \textbf{SEMESTRE:} \Semestre\par}
  \vspace{0.2cm}
  {\normalsize \textbf{PERIODO:} \Periodo\par}

  \vfill
  {\normalsize \LugarFecha\par}
\end{center}

\clearpage

% =========================
% ESTILO PARA EL RESTO DEL DOCUMENTO (Página 2 en adelante)
% =========================
\setcounter{page}{2}
\setlength{\headheight}{14.5pt} % evita warning fancyhdr
\pagestyle{fancy}
\fancyhf{}
\lhead{Residencia Profesional}
\rhead{Semana 1}
\cfoot{Página \thepage{} de \pageref{LastPage}}
\renewcommand{\headrulewidth}{0.4pt}

% =========================
% CONTENIDO
% =========================
\tableofcontents
\clearpage


\section{Proceso actual}

\subsection{Actores involucrados}
\begin{itemize}
  \item \textbf{Estudiantes} (nuevo ingreso o alumnos matriculados).
  \item \textbf{Personal de Servicio Social} (orientación, control de fila, captura de firma preliminar).
  \item \textbf{Operador de captura} (toma manual de fotografía).
  \item \textbf{Operadores de impresión} (genera PDFs, imprime, valida listas).
  \item \textbf{Servicios Escolares} (recibe listas y anuncia fechas de entrega).
\end{itemize}

\subsection{Flujo general}
\begin{enumerate}[leftmargin=*, label=\textbf{\arabic*.}]
  \item \textbf{Llegada y formación de filas.}
  \begin{itemize}
    \item Los estudiantes llegan y se forman; usualmente hay \textbf{filas largas}.
  \end{itemize}

  \item \textbf{Inducción manual de código de vestimenta.}
  \begin{itemize}
    \item El personal indica las reglas porque muchos estudiantes:
    \begin{itemize}
      \item no conocen el código de vestimenta,
      \item o no vieron/leyeron las indicaciones.
    \end{itemize}
    \item Indicaciones típicas:
    \begin{itemize}
      \item camisa tipo polo sin logos,
      \item cabello recogido/atrás de los hombros,
      \item sin lentes ni aretes,
      \item evitar prendas demasiado descubiertas
    \end{itemize}
  \end{itemize}

  \item \textbf{Recolección de firma mientras esperan.}
  \begin{itemize}
    \item Debido a la fila larga, se aprovecha para \textbf{recolectar firmas uno por uno}.
    \item Existe una \textbf{página web simple} para recolectar las firmas.
    \item Para ubicar al estudiante, se busca en \textbf{Discere} por:
    \begin{itemize}
      \item \textbf{número de control} o
      \item \textbf{número de ficha}.
    \end{itemize}
    \item Problemas observados:
    \begin{itemize}
      \item a veces la firma falla o no se registra bien,
      \item no hay un \textbf{mensaje claro de confirmación}.
    \end{itemize}
  \end{itemize}

  \item \textbf{Captura de foto.}
  \begin{itemize}
    \item Cuando el estudiante pasa a la estación de foto, se vuelve a pedir el identificador:
    \begin{itemize}
      \item ficha/control,
      \item o búsqueda por \textbf{nombre/carrera} (basado en opciones disponibles).
    \end{itemize}
    \item Un operador toma la foto manualmente.
    \item Problemas observados:
    \begin{itemize}
      \item errores al capturar o guardar la foto,
      \item si se repite la foto, se generan \textbf{múltiples imágenes} (una por intento).
    \end{itemize}
  \end{itemize}

  \item \textbf{Impresión de credenciales.}
  \begin{itemize}
    \item Desde el sistema se generan bloques de  \textbf{PDFs} y se aplica el formato de credencial correspondiente.
    \item Desde una computadora se manda a imprimir usando \textbf{parámetros específicos} para evitar atascos o errores de impresión.
    \item Se retrasa porque se valida con una \textbf{lista en papel} para confirmar qué credenciales ya se imprimieron.
    \item El proceso se repite por carrera, tiempo estimado de \textbf{3 meses}.
  \end{itemize}

  \item \textbf{Envío de listas a Servicios Escolares.}
  \begin{itemize}
    \item Se mandan listas a Servicios Escolares.
    \item Se anuncian fechas para la entrega de credenciales.
  \end{itemize}
\end{enumerate}

\subsection{Principales problemáticas del proceso actual}

\begin{itemize}
  \item \textbf{Saturación en la atención:} se presentan filas extensas, lo que incrementa los tiempos de espera y obliga a reiterar instrucciones de manera constante.
  \item \textbf{Duplicidad de actividades:} se realiza la identificación y/o búsqueda del estudiante en más de una etapa (captura de firma y, posteriormente, captura de fotografía), generando retrabajo.
  \item \textbf{Retroalimentación insuficiente del sistema:} no siempre se muestra una confirmación clara y verificable al momento de registrar la firma y/o la fotografía.
  \item \textbf{Gestión ineficiente de reintentos:} al repetir la toma de fotografía se generan múltiples archivos, sin garantizar de forma consistente que la evidencia final válida corresponda a la última captura.
  \item \textbf{Control manual de impresión y entrega:} el seguimiento del estado de impresión y entrega depende en gran medida de listas físicas, lo que incrementa el riesgo de errores y dificulta la trazabilidad.
\end{itemize}
\newpage
%==========================================================
\section{Proceso propuesto}

\subsection{Descripción general de la solución}
Se propone el desarrollo de un \textbf{sistema web institucional} (implementado en \textbf{PHP}) que permita digitalizar, ordenar y dar trazabilidad al proceso de credencialización, mediante los siguientes componentes:

\begin{itemize}
  \item \textbf{Sistema de turnos con código QR}, accesible desde teléfono móvil y/o visualizado en una pantalla pública.
  \item \textbf{Emisión de turno impreso} como alternativa para estudiantes que no cuenten con datos móviles.
  \item \textbf{Módulo de autoservicio (kiosco)}, para la confirmación de identidad y la aceptación de indicaciones previas a la captura.
  \item \textbf{Módulo de captura de evidencias} asociado al turno (firma, fotografía y huella).
  \item \textbf{Validación de calidad} previa a la impresión y \textbf{control de impresión/entrega} mediante registro en bitácora.
\end{itemize}

%==========================================================
\section{Flujo del proceso futuro por etapas}

\subsection{Etapa 1 --- Sistema de turnos con código QR (gestión de fila)}
\textbf{Propósito:} organizar la atención, reducir aglomeraciones y proporcionar visibilidad del avance del turno.

\textbf{Requerimientos funcionales:}
\begin{itemize}
  \item Generar un \textbf{turno de atención} mediante:
  \begin{itemize}
    \item código QR y portal web,
    \item acceso desde dispositivo móvil,
    \item y/o visualización en pantalla pública.
  \end{itemize}
  \item Mostrar información mínima del turno:
  \begin{itemize}
    \item turno actual en atención,
    \item tiempo estimado de espera.
  \end{itemize}
  \item Contemplar la \textbf{impresión de un turno} como alternativa para estudiantes sin conectividad.
\end{itemize}

\textbf{Reglas de operación propuestas:}
\begin{itemize}
  \item El turno tendrá \textbf{vigencia limitada.}
  \item Cada estudiante podrá generar más turnos durante el día; sin embargo, solo podrá contar con \textbf{un turno activo} a la vez.
  \item Si el turno vence, el estudiante podrá generar uno nuevo.
\end{itemize}

%==========================================================
\subsection{Etapa 2 --- Módulo de autoservicio (verificación previa)}
\textbf{Propósito:} disminuir errores y retrabajo, asegurando que el estudiante llegue a la captura con información confirmada.

\textbf{Requerimientos funcionales:}
\begin{itemize}
  \item Permitir que el estudiante se identifique mediante:
  \begin{itemize}
    \item número de control o número de ficha.
  \end{itemize}
  \item Visualizar y confirmar datos generales del estudiante y, en su caso:
  \begin{itemize}
    \item validar/actualizar información básica autorizada,
    \item presentar y solicitar aceptación de recomendaciones (código de vestimenta y guía para fotografía).
  \end{itemize}
  \item Permitir que este módulo pueda consultarse:
  \begin{itemize}
    \item en pantalla pública, o
    \item desde el dispositivo móvil del estudiante (cuando cuente con conectividad).
  \end{itemize}
\end{itemize}

%==========================================================
\subsection{Etapa 3 --- Módulo de captura de evidencias (firma, fotografía y huella)}
\textbf{Propósito:} capturar evidencias de manera controlada y vinculada a un turno vigente.

\textbf{Requerimientos funcionales:}
\begin{itemize}
  \item Capturar las evidencias requeridas:
  \begin{itemize}
    \item firma,
    \item fotografía,
    \item huella.
  \end{itemize}
  \item Asociar toda evidencia al \textbf{turno activo} (sesión por tiempo), evitando capturas aisladas sin control.
  \item Proporcionar retroalimentación clara:
  \begin{itemize}
    \item mensajes explícitos de éxito y error,
    \item control de reintentos,
    \item garantía de que la \textbf{evidencia válida final} corresponda a la última captura aprobada (especialmente en fotografía).
  \end{itemize}
\end{itemize}

%==========================================================
\subsection{Etapa 4 --- Validación de calidad (estudiante y personal)}
\textbf{Propósito:} garantizar la calidad mínima de la credencial antes de su impresión, reduciendo recapturas y reclamaciones.

\textbf{Requerimientos funcionales:}
\begin{itemize}
  \item Presentar una pantalla de validación posterior a la captura para revisar:
  \begin{itemize}
    \item fotografía (encuadre, nitidez, iluminación y visibilidad del rostro),
    \item firma,
    \item datos generales del estudiante.
  \end{itemize}
  \item Validación por el estudiante:
  \begin{itemize}
    \item confirmación de conformidad con datos y evidencias.
  \end{itemize}
  \item Validación por el personal:
  \begin{itemize}
    \item listado de credenciales pendientes de revisión,
    \item posibilidad de aprobar o enviar a recaptura cuando no cumpla criterios mínimos.
  \end{itemize}
\end{itemize}

%==========================================================
\subsection{Etapa 5 --- Control de impresión y respaldo en formato PDF}
\textbf{Propósito:} asegurar un flujo de impresión controlado y contar con alternativa operativa ante fallas.

\textbf{Requerimientos funcionales:}
\begin{itemize}
  \item Permitir la impresión posterior a la validación, con intervención mínima por parte del personal.
  \item Mantener un control de estados:
  \begin{itemize}
    \item credenciales pendientes por imprimir,
    \item credenciales impresas.
  \end{itemize}
  \item Proveer un mecanismo de respaldo:
  \begin{itemize}
    \item descarga manual del archivo PDF cuando la impresión automática no esté disponible.
  \end{itemize}
\end{itemize}

%==========================================================
\subsection{Etapa 6 --- Entrega, acuse y bitácora de seguimiento}
\textbf{Propósito:} registrar la entrega de credenciales con trazabilidad completa.

\textbf{Requerimientos funcionales:}
\begin{itemize}
  \item Registrar la entrega incluyendo:
  \begin{itemize}
    \item confirmación de entrega al estudiante,
    \item acuse (por definir: digital o físico),
    \item fecha y hora,
    \item responsable de la entrega,
    \item estado final de la credencial (entregada).
  \end{itemize}
  \item Permitir consulta para identificar:
  \begin{itemize}
    \item si la credencial ya fue impresa,
    \item si la credencial ya fue entregada.
  \end{itemize}
\end{itemize}

%==========================================================
\subsection{Etapa 7 --- Reportes e indicadores de operación}
\textbf{Propósito:} medir resultados, evidenciar mejoras y facilitar la toma de decisiones.

\textbf{Requerimientos sugeridos:}
\begin{itemize}
  \item Medición de tiempos por etapa:
  \begin{itemize}
    \item espera (turno),
    \item captura,
    \item validación,
    \item impresión,
    \item entrega.
  \end{itemize}
  \item Registro de reintentos:
  \begin{itemize}
    \item número de repeticiones de firma y fotografía.
  \end{itemize}
  \item Volumen operativo:
  \begin{itemize}
    \item estudiantes atendidos por día y por carrera.
  \end{itemize}
  \item Indicador de ``primera vez'':
  \begin{itemize}
    \item porcentaje de credenciales que no requieren recaptura ni reimpresión.
  \end{itemize}
\end{itemize}

%==========================================================
\section{Datos e identificación utilizados en el proceso}

\subsection{Datos generales del estudiante}
\begin{itemize}
  \item Nombre(s) y apellidos.
  \item Carrera.
  \item Firma.
  \item Fotografía.
  \item Número de control o número de ficha.
\end{itemize}

\subsection{Criterios de búsqueda e identificación}
\begin{itemize}
  \item Número de control.
  \item Número de ficha.
  \item Nombre.
  \item Carrera.
\end{itemize}

%==========================================================
\section{Situación operativa a resolver mediante el proceso futuro}
\begin{itemize}
  \item Reducir tiempos de espera y aglomeraciones mediante un control visible de turnos.
  \item Eliminar la duplicidad de identificación y captura en diferentes etapas del proceso.
  \item Asegurar confirmaciones claras y verificables al registrar firma y fotografía.
  \item Controlar reintentos y garantizar que la evidencia final válida sea la utilizada en la credencial.
  \item Fortalecer la trazabilidad de impresión y entrega mediante bitácora y estados de seguimiento.
\end{itemize}

%==========================================================
\section{Actividades realizadas (Semana 1)}
\begin{itemize}
  \item Levantamiento del proceso actual e identificación de problemáticas operativas.
  \item Documentación del proceso futuro por etapas como base para la definición de requerimientos y alcance.
  \item Identificación de datos de búsqueda y evidencias asociadas (firma, fotografía y huella).
\end{itemize}

\section{Pendientes para Semana 2}
\begin{itemize}
  \item Convertir las etapas del proceso futuro en requerimientos formales (historias de usuario y criterios de aceptación).
  \item Definir un modelo de datos mínimo (turno, estudiante, evidencias, estados y bitácora).
  \item Precisar el mecanismo de integración con el sistema institucional (Discere): campos disponibles, fuentes de datos y puntos de integración.
\end{itemize}
%==========================================================

\end{document}
